\documentclass[a4paper,11pt]{article}

% ============================================================
%  Le Collé — Séance 1 : Déterminants, courbes paramétrées
%  © 2026 Ahmed Amine Ouqlani
% ============================================================

% ---------- Polices ----------
\usepackage[utf8]{inputenc}
\usepackage[T1]{fontenc}
\usepackage[sc]{mathpazo}          % texte et mathématiques : Palatino
\usepackage[scaled=0.88]{helvet}   % titres et repères : linéale
\linespread{1.07}
\usepackage{amsmath,amssymb,amsthm,mathrsfs,stmaryrd}
\usepackage{microtype}

% ---------- Mise en page ----------
\usepackage[a4paper,top=2.3cm,bottom=2.3cm,left=2.4cm,right=2.4cm,
            headheight=14pt,headsep=0.8cm,footskip=1.1cm]{geometry}
\setlength{\parindent}{0pt}
\setlength{\parskip}{0.45em plus 0.15em minus 0.1em}
\raggedbottom
\emergencystretch=1.5em
\clubpenalty=10000 \widowpenalty=10000 \displaywidowpenalty=10000 \brokenpenalty=10000

% ---------- Couleurs : un bleu de nuit, un ocre, des gris ----------
\usepackage{xcolor}
\definecolor{Ink}{HTML}{1C2533}
\definecolor{Bleu}{HTML}{1F3A5F}
\definecolor{BleuMoyen}{HTML}{4A6A8F}
\definecolor{BleuClair}{HTML}{9DB2C9}
\definecolor{Fond}{HTML}{F2F4F7}
\definecolor{Ocre}{HTML}{B07A26}
\definecolor{Brique}{HTML}{A23B2A}
\definecolor{Gris}{HTML}{5E6672}
\definecolor{Filet}{HTML}{C5CCD6}

% ---------- Packages ----------
\usepackage{enumitem}
\setlist{itemsep=0.25em,topsep=0.35em,parsep=0pt}
\setlist[itemize]{label=\textcolor{BleuMoyen}{\small\textbullet},leftmargin=1.4em}
\setlist[enumerate]{leftmargin=1.8em,label=\textsf{\small\color{Bleu}\arabic*.}}
\usepackage{booktabs,array,multicol}
\usepackage{needspace,float}
\usepackage[font=small,labelfont={sf,color=Bleu},labelsep=period,skip=6pt]{caption}
\usepackage{tikz}
\usetikzlibrary{arrows.meta,calc,decorations.markings}
\usepackage{pgfplots}
\pgfplotsset{compat=1.17}
\usepackage{tkz-tab}
\usepackage[most]{tcolorbox}
\usepackage{titlesec,titletoc}
\usepackage{fancyhdr,lastpage,etoolbox}
\usepackage[colorlinks=true,linkcolor=Bleu,urlcolor=Bleu,citecolor=Bleu,
            pdftitle={Le Collé — TD Algèbre linéaire (corrections)},pdfauthor={Ahmed Amine Ouqlani}]{hyperref}

\numberwithin{equation}{section}
\numberwithin{figure}{section}
\numberwithin{table}{section}

% ---------- Titres ----------
% Chapitres (sections) : grand numéro, titre en linéale, filet ocre
\titleformat{\section}[display]
  {\normalfont\sffamily\color{Ink}}
  {{\small\color{Gris}\textls[200]{\MakeUppercase{Chapitre}}}\\[-0.2em]%
   {\fontsize{64}{64}\selectfont\color{BleuClair}\thesection}}
  {0.4em}
  {\fontsize{25}{30}\selectfont\bfseries}
  [\vspace{0.5em}{\color{Ocre}\rule{3.2cm}{1.4pt}}]
\titlespacing*{\section}{0pt}{0pt}{1.6em}
\titleformat{name=\section,numberless}[display]
  {\normalfont\sffamily\color{Ink}}{}{0pt}{\fontsize{25}{30}\selectfont\bfseries}
  [\vspace{0.5em}{\color{Ocre}\rule{3.2cm}{1.4pt}}]
\titlespacing*{name=\section,numberless}{0pt}{0pt}{1.6em}

\titleformat{\subsection}
  {\needspace{5\baselineskip}\normalfont\sffamily\Large\bfseries\color{Ink}}
  {\color{Bleu}\thesubsection}{0.8em}{}
\titlespacing*{\subsection}{0pt}{1.8em}{0.5em}
\titleformat{\subsubsection}
  {\needspace{5\baselineskip}\normalfont\sffamily\large\bfseries\color{Ink}}
  {\color{Bleu}\thesubsubsection}{0.7em}{}
\titlespacing*{\subsubsection}{0pt}{1.6em}{0.4em}
\titleformat{\paragraph}[runin]{\normalfont\sffamily\bfseries\color{Ink}}{}{0pt}{}[.]
\titlespacing*{\paragraph}{0pt}{0.9em}{0.6em}
\setcounter{secnumdepth}{3}
\setcounter{tocdepth}{2}

% ---------- Table des matières ----------
\titlecontents{section}[0pt]{\addvspace{0.7em}\sffamily\bfseries\color{Ink}}
  {\makebox[2.2em][l]{\color{Bleu}\thecontentslabel}}{}{\hfill\contentspage}[\addvspace{0.2em}]
\titlecontents{subsection}[2.2em]{\small\linespread{1}\selectfont}
  {\makebox[2.6em][l]{\color{Gris}\thecontentslabel}}{}{\titlerule*[0.7pc]{.}\contentspage}

% ---------- En-têtes et pieds de page ----------
\pagestyle{fancy}
\fancyhf{}
\renewcommand{\sectionmark}[1]{\markboth{\thesection\quad#1}{}}
\fancyhead[L]{\footnotesize\sffamily\color{Gris}\leftmark}
\fancyhead[R]{\footnotesize\sffamily\scshape\color{Gris}Le Collé}
\fancyfoot[C]{\small\sffamily\color{Ink}\thepage}
\fancyfoot[L]{\scriptsize\sffamily\hyperlink{toc}{\color{Filet}sommaire}}
\renewcommand{\headrulewidth}{0.4pt}
\renewcommand{\headrule}{{\color{Filet}\hrule height\headrulewidth width\headwidth}}
\renewcommand{\footrulewidth}{0pt}
\fancypagestyle{plain}{\fancyhf{}\fancyfoot[C]{\small\sffamily\color{Ink}\thepage}\renewcommand{\headrulewidth}{0pt}}

% ---------- Énoncés ----------
\newtheoremstyle{enonce}{0pt}{0pt}{}{}{\sffamily\bfseries\color{Bleu}}{.}{0.5em}
  {\thmname{#1}\thmnumber{~#2}\thmnote{\ \normalfont\sffamily\color{Ink}(#3)}}
\newtheoremstyle{exemple}{0.9em}{0.9em}{}{}{\sffamily\bfseries\color{BleuMoyen}}{.}{0.5em}
  {\thmname{#1}\thmnumber{~#2}\thmnote{\ \normalfont\sffamily\color{Ink}(#3)}}
\newtheoremstyle{remarque}{0.7em}{0.7em}{}{}{\sffamily\itshape}{.}{0.5em}{\thmname{#1}\thmnote{\ (#3)}}
\theoremstyle{enonce}
\newtheorem{definition}{Définition}[section]
\newtheorem{theoreme}[definition]{Théorème}
\newtheorem{proposition}[definition]{Proposition}
\newtheorem{corollaire}[definition]{Corollaire}
\theoremstyle{exemple}
\newtheorem{exemple}[definition]{Exemple}
\theoremstyle{remarque}
\newtheorem*{remarque}{Remarque}

\tcbset{
  defi/.style={enhanced,breakable,colback=white,colframe=white,boxrule=0pt,arc=0pt,
     borderline west={1.2pt}{0pt}{Bleu},left=11pt,right=4pt,top=3pt,bottom=3pt,
     before skip=1.1em,after skip=1.1em},
  thm/.style={enhanced,breakable,colback=Fond,colframe=Fond,boxrule=0pt,arc=0pt,
     borderline west={2.4pt}{0pt}{Bleu},left=11pt,right=8pt,top=6pt,bottom=6pt,
     before skip=1.2em,after skip=1.2em}}
\tcolorboxenvironment{definition}{defi}
\tcolorboxenvironment{theoreme}{thm}
\tcolorboxenvironment{proposition}{thm}
\tcolorboxenvironment{corollaire}{thm}

\renewcommand{\proofname}{Démonstration}
\renewcommand{\qedsymbol}{\textcolor{BleuClair}{$\blacksquare$}}
\makeatletter
\renewenvironment{proof}[1][\proofname]{\par\pushQED{\qed}\normalfont
  \topsep0.5em\relax\trivlist\item[\hskip\labelsep\sffamily\itshape\color{Gris}#1.]\ignorespaces}
  {\popQED\endtrivlist\@endpefalse}
\makeatother

% ---------- Méthode : filets haut et bas, titre en bandeau ----------
\makeatletter\newcommand{\methlabelfix}{\protected@edef\@currentlabel{\thetcbcounter}}\makeatother
\newtcolorbox[auto counter,number within=section]{methode}[2][]{%
  enhanced,unbreakable,colback=white,colframe=white,boxrule=0pt,arc=0pt,
  borderline north={0.9pt}{0pt}{Bleu},borderline south={0.9pt}{0pt}{Bleu},
  left=4pt,right=4pt,top=8pt,bottom=8pt,before skip=1.4em,after skip=1.4em,
  before upper={\methlabelfix\phantomsection\addcontentsline{met}{metentry}{\protect\makebox[2.6em][l]{\thetcbcounter}#2}%
     {\sffamily\bfseries\small\color{white}\colorbox{Bleu}{\strut\ MÉTHODE~\thetcbcounter\ }}%
     \ \ {\sffamily\bfseries\color{Bleu}#2}\par\smallskip},#1}

% ---------- Repères pédagogiques : un filet de couleur, un mot-clé ----------
\newcommand{\definirbloc}[3]{% nom, couleur, libellé
  \newtcolorbox{#1}[1][]{enhanced,breakable,blanker,
    borderline west={1.4pt}{0pt}{#2},left=10pt,right=2pt,top=1pt,bottom=1pt,
    before skip=1em,after skip=1em,
    before upper={{\sffamily\bfseries\footnotesize\color{#2}\textls[120]{\MakeUppercase{#3}}}\quad},##1}}
\definirbloc{reflexe}{Bleu}{Réflexe}
\definirbloc{astuce}{Ocre}{Astuce}
\definirbloc{attention}{Brique}{Attention}
\definirbloc{enpratique}{BleuMoyen}{En pratique}
\definirbloc{plusloin}{Gris}{Pour aller plus loin}
\definirbloc{strategie}{Bleu}{Stratégie}
\definirbloc{lecture}{BleuMoyen}{Lire le schéma}
% compatibilité avec les parties reprises
\newenvironment{garde}{\begin{attention}}{\end{attention}}
\newenvironment{complementHP}{\begin{plusloin}}{\end{plusloin}}
\newenvironment{interpretation}{\begin{enpratique}}{\end{enpratique}}

% ---------- Exercices, corrigés ----------
\newcounter{exo}[section]
\renewcommand{\theexo}{\thesection.\arabic{exo}}
\newcommand{\niveau}[1]{{\color{Ocre}\footnotesize
  \ifcase#1\or$\bullet\circ\circ\circ\circ$\or$\bullet\bullet\circ\circ\circ$\or$\bullet\bullet\bullet\circ\circ$%
  \or$\bullet\bullet\bullet\bullet\circ$\or$\bullet\bullet\bullet\bullet\bullet$\fi}}
\newcommand{\serie}[1]{\par\needspace{8\baselineskip}\vspace{1.5em}%
  \noindent{\sffamily\bfseries\color{Bleu}#1}\par\nobreak\vspace{-0.4em}%
  \noindent{\color{Filet}\rule{\linewidth}{0.4pt}}\par\nobreak\vspace{0.1em}}
\newenvironment{exo}[3]{%
  \par\needspace{5\baselineskip}\vspace{0.75em}\refstepcounter{exo}\label{e:#1}%
  \def\exolab{#1}%
  \noindent{\sffamily\bfseries\color{Bleu}Exercice~\theexo}\quad\niveau{#2}%
  \hfill{\footnotesize\sffamily\color{Gris}#3}\par\nobreak\vspace{0.15em}}%
  {\par\nopagebreak\noindent\hfill{\scriptsize\sffamily\hyperref[c:\exolab]{\color{Gris}corrigé p.~\pageref{c:\exolab}}}\par}
\newenvironment{corrige}[1]{%
  \par\needspace{6\baselineskip}\vspace{1em}\phantomsection\label{c:#1}%
  \noindent{\sffamily\bfseries\color{Bleu}Exercice~\ref{e:#1}}\quad
  {\scriptsize\sffamily\hyperref[e:#1]{\color{Gris}énoncé p.~\pageref{e:#1}}}\par\nobreak\vspace{0.2em}}%
  {\hfill\textcolor{BleuClair}{$\square$}\par}
\newcommand{\etape}[1]{\par\smallskip\noindent{\sffamily\bfseries\small\color{BleuMoyen}#1.}\ }
\newenvironment{solution}{\par\noindent{\sffamily\bfseries\small\color{BleuMoyen}Solution.}\ }{\hfill\textcolor{BleuClair}{$\square$}\par}
\newcommand{\chercher}{\par\smallskip\noindent{\color{Filet}\rule[0.5ex]{1.2cm}{0.4pt}}\
  {\footnotesize\sffamily\itshape\color{Gris}laisser chercher l'étudiant avant de poursuivre}\par\smallskip}

% ---------- Guide de séance ----------

\newcommand{\marge}[1]{\leavevmode\marginpar{\raggedleft\scriptsize\sffamily\color{Gris}\textls[100]{\MakeUppercase{#1}}}\ignorespaces}
\newcommand{\Etape}[4]{%
  \subsection[Étape #1 --- #3 {\color{Gris}(#2)}]{#3}\label{etape:#1}%
  \noindent{\small\sffamily\color{Gris}{\color{Ocre}\bfseries Étape #1}\quad·\quad #2\quad·\quad #4}\par\medskip}
\newcommand{\question}[1]{\par\smallskip\noindent{\sffamily\bfseries\small\color{BleuMoyen}Question.}\ \textit{#1}\par\smallskip}
\newcommand{\rappel}[1]{\par\smallskip\noindent\phantomsection\addcontentsline{rap}{rapentry}{#1}%
  {\sffamily\bfseries\scriptsize\colorbox{Fond}{\color{Gris}\strut RAPPEL}}\ {\sffamily\small\color{Gris}#1.}\ }
\newcommand{\pourquoi}{\par\noindent{\sffamily\bfseries\small\color{Ocre}Pourquoi ?}\ }

% Ouverture de chapitre : \ouverture{objectifs}{prérequis}{durée}
\newcommand{\ouverture}[3]{%
  \begin{tcolorbox}[enhanced,colback=Fond,colframe=Fond,arc=0pt,boxrule=0pt,left=12pt,right=12pt,top=10pt,bottom=10pt,
     before skip=1.2em,after skip=1.6em]
  \begin{minipage}[t]{0.56\linewidth}
  {\sffamily\bfseries\small\color{Bleu}OBJECTIFS}\par\vspace{0.3em}\small #1
  \end{minipage}\hfill
  \begin{minipage}[t]{0.38\linewidth}
  {\sffamily\bfseries\small\color{Bleu}PRÉREQUIS}\par\vspace{0.3em}\small #2
  \ifx\relax#3\relax\else\par\vspace{0.6em}{\sffamily\bfseries\small\color{Bleu}DURÉE}\par\vspace{0.2em}\small #3\fi
  \end{minipage}
  \end{tcolorbox}}
% Fin de chapitre
\newenvironment{essentiel}{\par\needspace{12\baselineskip}\vspace{1.5em}%
  \begin{tcolorbox}[enhanced,unbreakable,colback=white,colframe=white,boxrule=0pt,arc=0pt,
   borderline north={1.4pt}{0pt}{Ocre},left=2pt,right=2pt,top=8pt,bottom=4pt]
  {\sffamily\bfseries\color{Ocre}L'ESSENTIEL DU CHAPITRE}\par\medskip}{\end{tcolorbox}}

% ---------- Notations ----------
\newcommand{\K}{\mathbb{K}}
\newcommand{\R}{\mathbb{R}}
\newcommand{\C}{\mathbb{C}}
\newcommand{\N}{\mathbb{N}}
\newcommand{\Z}{\mathbb{Z}}
\newcommand{\Mn}[1]{\mathcal{M}_{n}(#1)}
\newcommand{\GLn}[1]{\mathrm{GL}_n(#1)}
\newcommand{\LE}{\mathcal{L}(E)}
\newcommand{\rg}{\operatorname{rg}}
\newcommand{\vect}{\operatorname{Vect}}
\newcommand{\Mat}{\operatorname{Mat}}
\newcommand{\Id}{\operatorname{Id}}
\newcommand{\sh}{\operatorname{sh}}
\newcommand{\ch}{\operatorname{ch}}
\renewcommand{\th}{\operatorname{th}}
\DeclareMathOperator{\argsh}{argsh}
\newcommand{\vv}[1]{\overrightarrow{#1}}
\newcommand{\ve}[1]{\vec{#1}}
\newcommand{\ient}[2]{\llbracket #1,#2\rrbracket}
\newcommand{\dt}[1]{\left|\begin{matrix}#1\end{matrix}\right|}
\newcommand{\dis}{\displaystyle}

% ---------- Figures ----------
\pgfplotsset{trig format plots=rad}
\pgfplotsset{courbe/.style={axis lines=middle,axis equal image,
  axis line style={Gris,-{Stealth[length=5pt]}},
  xlabel={$x$},ylabel={$y$},ticklabel style={font=\scriptsize\sffamily,color=Gris},
  every axis x label/.style={at={(ticklabel* cs:1)},anchor=west,font=\small},
  every axis y label/.style={at={(ticklabel* cs:1)},anchor=south,font=\small},
  samples=300,smooth,width=0.6\linewidth}}
\tikzset{arcetude/.style={Bleu,line width=1.25pt},
         reste/.style={BleuClair,line width=1pt},
         tgte/.style={Ocre,line width=0.9pt},
         asy/.style={Gris,dashed,line width=0.6pt},
         pt/.style={circle,fill=Ink,inner sep=1.2pt}}

% ---------- Navigation ----------
\makeatletter
\newcommand*\l@rapentry{\@dottedtocline{1}{0em}{0em}}
\newcommand*\l@metentry{\@dottedtocline{1}{0em}{0em}}
\newcommand{\listemethodes}{\@starttoc{met}}
\newcommand{\listerappels}{\@starttoc{rap}}
\makeatother
% onglets de chapitre sur le bord droit
\usepackage{eso-pic}
\newif\ifonglets \ongletsfalse
\newcommand{\dessineronglet}{%
  \edef\ongcol{\ifnum\value{section}=3 Ocre\else Bleu\fi}%
  \edef\ongnum{\the\value{section}}%
  \pgfmathsetmacro{\ongy}{3.2+2.3*(\ongnum-1)}%
  \begin{tikzpicture}[remember picture,overlay]
  \fill[\ongcol] ([xshift=-0.75cm,yshift=-\ongy cm]current page.north east)
     rectangle ([yshift=-\ongy cm-1.9cm]current page.north east);
  \node[text=white,font=\sffamily\bfseries\large] at ([xshift=-0.37cm,yshift=-\ongy cm-0.95cm]current page.north east)
     {\ongnum};
  \end{tikzpicture}}
\AddToShipoutPictureBG{\ifonglets\ifnum\value{section}>0 \dessineronglet\fi\fi}
% blocs de navigation et d'enseignement
\definirbloc{commencer}{Gris}{Si vous commencez ici}
\definirbloc{horsprogramme}{Brique}{Hors programme}
\newcommand{\HP}{\,{\sffamily\bfseries\scriptsize\textcolor{Brique}{[HP]}}}
\newcommand{\source}[1]{{\footnotesize\sffamily\color{Gris}\,(#1)}}
\newcommand{\coupdepouce}[1]{\par\vspace{-0.2em}\noindent{\footnotesize\sffamily\color{Gris}%
  \textbf{Coup de pouce} (à donner si l'étudiant bloque) : #1}\par}
\newcommand{\pointarret}[1]{\par\vspace{1em}\noindent{\color{Filet}\rule[0.5ex]{\linewidth}{0.4pt}}\par\nobreak
  \noindent{\footnotesize\sffamily\color{Ocre}\bfseries POINT D'ARRÊT POSSIBLE}\quad{\footnotesize\sffamily\color{Gris}#1}\par\vspace{0.6em}}
\newcommand{\demarche}[1]{\noindent{\sffamily\bfseries\small\color{BleuMoyen}Démarche.}\ \textit{#1}\par\smallskip}

% ---------- Corrigés : idées, astuces, pièges, à retenir ----------
\newcommand{\choix}[1]{\par\noindent{\sffamily\bfseries\small\color{BleuMoyen}Choix de la méthode.}\ #1\par\smallskip}
\newcommand{\gain}[1]{\par\noindent{\sffamily\bfseries\small\color{Ocre}Astuce.}\ #1\par}
\newcommand{\piege}[1]{\par\noindent{\sffamily\bfseries\small\color{Brique}Piège.}\ #1\par}
\newtcolorbox{retenirbox}{enhanced,breakable,colback=Fond,colframe=Fond,arc=0pt,boxrule=0pt,
  borderline west={2pt}{0pt}{Ocre},left=8pt,right=6pt,top=4pt,bottom=4pt,before skip=0.7em,after skip=0.5em}
\newcommand{\retenir}[1]{\begin{retenirbox}{\sffamily\bfseries\footnotesize\color{Ocre}À RETENIR}\quad #1\end{retenirbox}}
\newcommand{\TD}[1]{{\color{Ocre}\bfseries TD\,#1}\ ·\ }
\newcommand{\q}[1]{\par\smallskip\noindent{\sffamily\bfseries\color{Bleu}#1}\ }
\providecommand{\iddots}{\mathinner{\mkern1mu\raise1pt\hbox{.}\mkern2mu\raise4pt\hbox{.}\mkern2mu\raise7pt\hbox{.}\mkern1mu}}

\providecommand{\Ker}{\operatorname{Ker}}\renewcommand{\Im}{\operatorname{Im}}\providecommand{\tr}{\operatorname{tr}}\providecommand{\Sp}{\operatorname{Sp}}

% ================================================================ TD : mise en page propre au livret d'exercices
\fancyhf{}
\fancyhead[L]{\footnotesize\sffamily\color{Gris}TD · Algèbre linéaire · corrections}
\fancyhead[R]{\footnotesize\sffamily\scshape\color{Gris}Le Collé}
\fancyfoot[C]{\small\sffamily\color{Ink}\thepage\,/\,\pageref{LastPage}}
\fancypagestyle{plain}{\fancyhf{}\fancyfoot[C]{\small\sffamily\color{Ink}\thepage\,/\,\pageref{LastPage}}\renewcommand{\headrulewidth}{0pt}}
\makeatletter
\newcommand{\partie}[3]{% numéro, titre, description
  \par\needspace{10\baselineskip}\vspace{1.4em}\phantomsection\label{partie:#1}\pdfbookmark[0]{Partie #1}{partie.#1}%
  {\sffamily\small\color{Gris}\textls[200]{PARTIE #1}}\par\nobreak\vspace{0.1em}
  {\sffamily\bfseries\LARGE\color{Ink}#2}\par\nobreak\vspace{0.3em}
  {\color{Ocre}\rule{2.6cm}{1.2pt}}\par\nobreak\vspace{0.2em}
  {\small\color{Gris}#3}\par\nobreak\vspace{0.6em}}
\newcommand{\exercice}[7]{% numéro, difficulté, titre, origine, notions, identifiant, page du corrigé
  \par\needspace{8\baselineskip}\vspace{1.1em}\phantomsection
  \def\@currentlabel{#1}\label{e:#6}\pdfbookmark[1]{Exercice #1}{e.#6}%
  \noindent{\sffamily\bfseries\large\color{Bleu}Exercice~#1}\quad\niveau{#2}\hfill{\footnotesize\sffamily\color{Gris}#4}\par\nobreak
  \vspace{0.1em}\noindent{\sffamily\bfseries\color{Ink}#3}\par\nobreak
  \vspace{0.1em}\noindent{\footnotesize\sffamily\color{Gris}Notions : #5.\hfill\mbox{corrigé p.~#7}\par}\nobreak
  \vspace{0.35em}}
\newcommand{\renvoi}[1]{\par\nopagebreak\vspace{0.2em}\noindent\hfill{\scriptsize\sffamily\color{Gris}#1}\par}
\newcommand{\place}{\par\vspace{1.6cm plus 0.6cm minus 0.6cm}}
\newcommand{\correction}[4]{% numéro, titre, identifiant, renvoi
  \par\needspace{7\baselineskip}\vspace{1.3em}\phantomsection
  \def\@currentlabel{#1}\label{e:#3}\pdfbookmark[1]{Exercice #1}{c.#3}%
  {\color{Filet}\hrule height 0.4pt}\vspace{0.5em}
  \noindent{\sffamily\bfseries\large\color{Bleu}Exercice~#1}\quad{\sffamily\bfseries\color{Ink}#2}\hfill{\scriptsize\sffamily\color{Gris}#4}\par\nobreak
  \vspace{0.4em}}
\newcommand{\sommaireligne}[3]{\noindent\hbox to 2.2em{\hfill\sffamily\color{Bleu}#1.}\enspace\begin{minipage}[t]{0.8\linewidth}\raggedright#2\end{minipage}\hfill\pageref{e:#3}\par}
\makeatother
\setlength{\parskip}{0.5em plus 0.15em minus 0.1em}

\begin{document}
\thispagestyle{plain}
{\sffamily\small\color{Gris}\textls[150]{LE COLLÉ · MATHÉMATIQUES PT · CHAPITRE 1}}\par\vspace{0.2em}
{\sffamily\bfseries\fontsize{26}{30}\selectfont\color{Ink}TD --- Algèbre linéaire}\par\vspace{0.3em}
{\sffamily\Large\color{BleuMoyen}Corrections}\par\vspace{0.4em}
{\color{Ocre}\rule{3.2cm}{1.4pt}}\par\vspace{0.8em}
{\small Corrections détaillées des 19 exercices du TD, dans le même ordre et avec les mêmes numéros.
Chaque correction commence par le choix de la méthode, puis détaille les calculs ; les encadrés « À retenir » résument ce qui
resservira. Les renvois « méthode~1.3 du livret » désignent le cours.
\textbf{Chercher d'abord l'exercice} avant d'ouvrir sa correction.}

\partie{I}{Applications directes du cours}{}
\correction{1}{Un hyperplan de $\mathbb K^n$ et une base}{aln-1}{énoncé : TD, p.~2}
\choix{$H$ est le noyau d'une forme linéaire non nulle : c'est un hyperplan ; il suffit ensuite d'une famille libre de bon cardinal.}
\q{1.} $\varphi:(x_i)\mapsto\sum x_i$ est une forme linéaire non nulle ($\varphi(e_1)=1$) : $H=\Ker\varphi$ est un hyperplan, $\dim H=n-1$.
\q{2.} Chaque $f_i$ est dans $H$ ($1-1=0$). Si $\sum_{i=1}^{n-1}\lambda_if_i=0$, la $i$-ième coordonnée vaut $\lambda_i=0$ pour $i\leq n-1$ :
la famille est libre. Elle a $n-1=\dim H$ vecteurs : c'est une base de $H$.
\retenir{famille libre de $\dim H$ vecteurs de $H$ : base de $H$, sans vérifier qu'elle engendre.}
\par\hfill\textcolor{BleuClair}{$\square$}\par
\correction{2}{Un supplémentaire dans $\mathbb R_2[X]$}{aln-2}{énoncé : TD, p.~2}
$P\mapsto P(1)$ est une forme linéaire non nulle sur $\R_2[X]$ : $F$ est un hyperplan, de dimension $2$. Le polynôme constant $1$
n'est pas dans $F$, donc $\R_2[X]=F\oplus\vect(1)$ (un hyperplan et une droite qui n'y est pas contenue). Décomposition :
$P=P(1)+\big(P-P(1)\big)$, avec $P(1)\in\vect(1)$ et $P-P(1)\in F$.
\retenir{hyperplan $+$ droite non contenue dans l'hyperplan $=$ supplémentaires.}
\par\hfill\textcolor{BleuClair}{$\square$}\par
\correction{3}{Sommes directes dans $\mathbb R^3$}{red-somme}{énoncé : TD, p.~2}
\q{1.} Si $v\in F\cap G$, $v=(m,m,m)$ et $3m=0$ : $F\cap G=\{0\}$. De plus $\dim F=2$ (un hyperplan de $\R^3$) et $\dim G=1$ : la
somme est directe et de dimension $3$, donc $\R^3=F\oplus G$. On cherche $v=w+m(1,1,1)$ avec $w\in F$ : la somme des coordonnées
de $w$ est nulle, d'où $x+y+z=3m$. Donc
\[v=\underbrace{\big(x-m,\,y-m,\,z-m\big)}_{\in F}+\underbrace{m(1,1,1)}_{\in G},\qquad m=\frac{x+y+z}3.\]
\q{2.} Les intersections deux à deux sont nulles, mais $(1,0,0)+(0,1,0)-(1,1,0)=0$ est une décomposition non triviale de $0$ :
la somme n'est pas directe (proposition~1.7 du livret). D'ailleurs $D_1+D_2+D_3$ est le plan $z=0$, de dimension $2<3$.
\par\hfill\textcolor{BleuClair}{$\square$}\par
\correction{4}{Hyperplans de $\mathbb R^4$}{red-hyp}{énoncé : TD, p.~3}
\q{1.} En ajoutant et en retranchant les équations : $x+z=0$ et $y+t=0$. Donc $H_1\cap H_2=\{(x,y,-x,-y)\}$, de base
$\big((1,0,-1,0),(0,1,0,-1)\big)$ : dimension $2=4-2$.
\q{2.} $u\notin H_1$ ($1+1+1+1=4\neq0$) : $\R^4=H_1\oplus\vect(u)$ (proposition~1.11 du livret). En revanche $u\in H_2$
($1-1+1-1=0$) : $\vect(u)\subset H_2$, ce n'est pas un supplémentaire.
\par\hfill\textcolor{BleuClair}{$\square$}\par
\correction{5}{Projecteurs associés}{aln-12}{énoncé : TD, p.~3}
Pour $x=x_F+x_G$ (décomposition unique), $p(x)=x_F$ et $q(x)=x_G$ : $p(x)+q(x)=x$. Ensuite $q(x)\in G$, donc $p\big(q(x)\big)=0$ ;
de même $q\circ p=0$. Enfin $\Ker p=G=\Im q$ et $\Ker q=F=\Im p$.
\retenir{une décomposition $E=F\oplus G$ donne deux projecteurs de somme $\Id_E$.}
\par\hfill\textcolor{BleuClair}{$\square$}\par
\correction{6}{Fonctions paires et impaires}{rdn-17}{énoncé : TD, p.~3}
\choix{dimension infinie : on revient à la définition (existence et unicité de la décomposition), par analyse-synthèse.}
\emph{Unicité.} Si $f\in\mathcal P\cap\mathcal I$, $f(-x)=f(x)=-f(x)$, donc $f=0$. \emph{Existence.} Pour $f\in E$, on pose
$p(x)=\frac{f(x)+f(-x)}2$ (paire) et $i(x)=\frac{f(x)-f(-x)}2$ (impaire) : $f=p+i$. Donc $E=\mathcal P\oplus\mathcal I$.
Pour $\exp$ : $\exp=\ch+\sh$.
\gain{l'application $s:f\mapsto\big(x\mapsto f(-x)\big)$ vérifie $s\circ s=\Id_E$ : c'est une symétrie, et $\mathcal P=\Ker(s-\Id_E)$,
$\mathcal I=\Ker(s+\Id_E)$. La décomposition est celle de la proposition sur les symétries.}
\par\hfill\textcolor{BleuClair}{$\square$}\par
\partie{II}{Exercices classiques et méthodes incontournables}{}
\correction{7}{Projecteur et symétrie dans $\mathbb R^3$}{al-proj}{énoncé : TD, p.~4}
\choix{on décompose un vecteur quelconque selon $P\oplus D$ : la condition « $x-x_D\in P$ » donne directement la composante sur $D$.}
\q{1.} $P$ : $y=2x+z$, de base $\big((1,2,0),(0,1,1)\big)$. $D$ : $(x,0,x)$, de base $w=(1,0,1)$. Comme $2\cdot1-0+1=3\neq0$,
$w\notin P$ : $P\cap D=\{0\}$, et $\dim P+\dim D=3$, donc $\R^3=P\oplus D$ (c'est aussi la proposition~1.11 du livret, point 2).
\q{2.} Soit $v=(x,y,z)$. On cherche $t$ tel que $v-tw\in P$ : $2(x-t)-y+(z-t)=0$, soit $t=\frac{2x-y+z}3$. Alors
$p(v)=v-tw$ et $s(v)=2p(v)-v=v-2tw$. Dans la base canonique :
\[\Mat(p)=\frac13\begin{pmatrix}1&1&-1\\0&3&0\\-2&1&2\end{pmatrix},\qquad\Mat(s)=\frac13\begin{pmatrix}-1&2&-2\\0&3&0\\-4&2&1\end{pmatrix}.\]
Le calcul direct donne $\Mat(p)^2=\Mat(p)$ et $\Mat(s)^2=I_3$. Contrôles : $\tr p=2=\rg p=\dim P$, et $\tr s=2-1=1$.
\retenir{pour une projection, on ne cherche que la composante « parallèle » ; la symétrie s'en déduit par $s=2p-\Id$.}
\par\hfill\textcolor{BleuClair}{$\square$}\par
\correction{8}{Deux propriétés de la trace}{red-trace}{énoncé : TD, p.~4}
\q{1.} Par linéarité et $\tr(AB)=\tr(BA)$ : $\tr(AB-BA)=0$, alors que $\tr(I_n)=n\neq0$.
\q{2.} L'endomorphisme $p$ de $\K^n$ associé à $A$ est un projecteur ; dans une base adaptée à $\K^n=\Im p\oplus\Ker p$, sa matrice
est $\mathrm{diag}(I_r,0)$ avec $r=\rg(p)$. Donc $A$ est semblable à cette matrice, et $\tr A=r=\rg A$.
\retenir{la trace ne voit pas l'ordre d'un produit de deux matrices : $\tr(AB-BA)=0$.}
\par\hfill\textcolor{BleuClair}{$\square$}\par
\correction{9}{La trace détermine une matrice}{aln-5}{énoncé : TD, p.~4}
\choix{tester la relation sur les matrices de la base canonique.}
Le sens réciproque est clair. Pour $M=E_{j,i}$, la matrice $AE_{j,i}$ a pour $i$-ième colonne la $j$-ième colonne de $A$ et des
colonnes nulles ailleurs : son seul coefficient diagonal éventuellement non nul est en position $(i,i)$, et vaut $a_{i,j}$. Donc
$\tr(AE_{j,i})=a_{i,j}$, et l'hypothèse donne $a_{i,j}=b_{i,j}$ pour tous $i,j$.
\retenir{$\tr(AE_{j,i})=a_{i,j}$ : les formes $M\mapsto\tr(AM)$ suffisent à retrouver $A$.}
\par\hfill\textcolor{BleuClair}{$\square$}\par
\correction{10}{Rang d'une matrice à paramètres}{aln-6}{énoncé : TD, p.~5}
\choix{opérations élémentaires sur les lignes et les colonnes, qui ne changent pas le rang ; on discute aux valeurs qui annulent un pivot.}
$C_2\leftarrow C_2-bC_3$ donne les colonnes $(a,1,1)$, $(0,ab-b,b-ab)=b(a-1)(0,1,-1)$ et $(1,1,a)$.
\begin{itemize}
\item Si $a=1$ : les trois lignes de $M$ sont égales à $(1,b,1)$ : $\rg M=1$.
\item Si $a\neq1$ et $b=0$ : la deuxième colonne est nulle, et $(a,1,1)$, $(1,1,a)$ ne sont pas proportionnelles ($a\neq1$) : $\rg M=2$.
\item Si $a\neq1$ et $b\neq0$ : $\rg M$ est le rang de la matrice de lignes $(a,0,1)$, $(1,1,1)$, $(1,-1,a)$. On prend $(1,1,1)$
comme pivot : $L_1\leftarrow L_1-aL_2$ et $L_3\leftarrow L_3-L_2$ donnent $(0,-a,1-a)$ et $(0,-2,a-1)$. Ces deux lignes sont
proportionnelles si et seulement si $-a(a-1)+2(1-a)=-(a-1)(a+2)=0$, soit $a=-2$. Donc $\rg M=2$ si $a=-2$, et $\rg M=3$ sinon.
\end{itemize}
\gain{au chapitre~2, on obtiendra d'un coup $\det M=b(a-1)^2(a+2)$, qui donne les cas de rang $3$.}
\par\hfill\textcolor{BleuClair}{$\square$}\par
\correction{11}{Un projecteur donné par sa matrice}{aln-4}{énoncé : TD, p.~5}
\q{1.} Le calcul donne $A^2=A$ : $u$ est un projecteur, et $\rg u=\tr A=\frac12+1+\frac12=2$.
\q{2.} $\Ker u$ : $\frac x2+y-\frac z2=0$, $y=0$, d'où $x=z$ : $\Ker u=\vect(1,0,1)$. $\Im u$ est engendrée par les colonnes, par
exemple $(1,0,-1)$ (le double de la première) et $(1,1,1)$, non proportionnelles : c'est une base. Un vecteur $(a,b,c)$ orthogonal
aux deux, $(1,-2,1)$, donne l'équation $x-2y+z=0$ (les deux vecteurs la vérifient, et $\dim\Im u=2$).
\q{3.} Pour un projecteur, $\R^3=\Ker u\oplus\Im u$ et $u$ vaut $0$ sur $\Ker u$, $\Id$ sur $\Im u$. Dans la base adaptée
$\big((1,0,1),(1,0,-1),(1,1,1)\big)$, la matrice de $u$ est $\mathrm{diag}(0,1,1)$ : $P=\begin{pmatrix}1&1&1\\0&0&1\\1&-1&1\end{pmatrix}$.
\retenir{projecteur : rang $=$ trace, et la base adaptée à $\Ker\oplus\Im$ donne la matrice $\mathrm{diag}(0,\dots,0,1,\dots,1)$.}
\par\hfill\textcolor{BleuClair}{$\square$}\par
\correction{12}{Lire des blocs}{red-stable}{énoncé : TD, p.~5}
\q{1.} $f(x,0,z)=(3x+z,\,0,\,x+3z)$ a une deuxième coordonnée nulle : $F$ est stable. $f(u)=A(1,1,0)^T=(2,2,0)=2u$ : $G$ est stable.
\q{2.} $u\notin F$ (sa deuxième coordonnée vaut $1$) : $\R^3=F\oplus G$. Dans la base $(e_1,e_3,u)$ : $f(e_1)=(3,0,1)=3e_1+e_3$,
$f(e_3)=(1,0,3)=e_1+3e_3$, $f(u)=2u$, donc
\[\Mat(f)=\begin{pmatrix}3&1&0\\1&3&0\\0&0&2\end{pmatrix},\qquad\tr A=3+3+2=8.\]
\retenir{deux sous-espaces stables supplémentaires $\Rightarrow$ matrice diagonale par blocs.}
\par\hfill\textcolor{BleuClair}{$\square$}\par
\partie{III}{Exercices de synthèse et approfondissement}{}
\correction{13}{Deux sous-espaces trop grands se coupent}{aln-3}{énoncé : TD, p.~6}
Par la formule de Grassmann, $\dim(F\cap G)=\dim F+\dim G-\dim(F+G)\geq\dim F+\dim G-\dim E>0$. Deux plans de $\R^3$ :
$2+2>3$, donc leur intersection est de dimension au moins $1$.
\par\hfill\textcolor{BleuClair}{$\square$}\par
\correction{14}{Matrices de rang au plus $1$}{aln-7}{énoncé : TD, p.~6}
Si $M=UV^T$, la $j$-ième colonne de $M$ est $v_jU$ : toutes les colonnes sont dans $\vect(U)$, et $\rg M\leq1$. Réciproquement, si
$\rg M=0$, $M=0=0\cdot V^T$ ; si $\rg M=1$, l'image est une droite $\vect(U)$, chaque colonne s'écrit $v_jU$, et $M=UV^T$ avec
$V=(v_j)$. Enfin $V^TU=\sum_iv_iu_i=\tr(UV^T)$ est un scalaire, donc $M^2=U(V^TU)V^T=\tr(M)\,M$.
\retenir{une matrice de rang $1$ est un produit « colonne $\times$ ligne », et $M^2=\tr(M)\,M$.}
\par\hfill\textcolor{BleuClair}{$\square$}\par
\correction{15}{Une matrice de carré nul}{aln-8}{énoncé : TD, p.~6}
\q{1.} Les trois lignes sont proportionnelles à $(1,1,-1)$ : $\rg f=1$, et $\Ker f$ est le plan $x+y-z=0$, de base
$\big((-1,1,0),(1,0,1)\big)$. $\Im f=\vect(u)$, $u=(1,-3,-2)$. Comme $1-3+2=0$, $u\in\Ker f$ : $\Im f\subset\Ker f$, donc $f\circ f=0$.
\q{2.} $f(e_1)=u$. Dans la base $(u,k,e_1)$ avec $k=(1,0,1)$ : $(u,k)$ est une base de $\Ker f$ et $e_1\notin\Ker f$, c'est donc une
base de $\R^3$. Comme $f(u)=f(k)=0$ et $f(e_1)=u$, la matrice de $f$ a pour seul coefficient non nul un $1$ en position $(1,3)$.
\retenir{$\Ker f$ et $\Im f$ ne sont pas toujours supplémentaires : ici $\Im f\subset\Ker f$.}
\par\hfill\textcolor{BleuClair}{$\square$}\par
\correction{16}{L'endomorphisme $M\mapsto AM-MA$}{al-crochet}{énoncé : TD, p.~7}
\choix{on calcule l'image de chaque vecteur de la base canonique ; la trace donne le contrôle de l'image.}
\q{1.} $\varphi(\lambda M+N)=A(\lambda M+N)-(\lambda M+N)A=\lambda\varphi(M)+\varphi(N)$, et $\varphi(M)\in\mathcal M_2(\R)$. Par
$\tr(AM)=\tr(MA)$, $\tr\varphi(M)=0$ : $\Im\varphi$ est contenue dans l'hyperplan $\Ker\tr$, de dimension $3$.
\q{2.} $\varphi(E_{1,1})=\begin{pmatrix}0&-1\\1&0\end{pmatrix}$, $\varphi(E_{1,2})=\begin{pmatrix}-1&2\\0&1\end{pmatrix}$,
$\varphi(E_{2,1})=\begin{pmatrix}1&0\\-2&-1\end{pmatrix}$, $\varphi(E_{2,2})=\begin{pmatrix}0&1\\-1&0\end{pmatrix}$, d'où
\[\Mat(\varphi)=\begin{pmatrix}0&-1&1&0\\-1&2&0&1\\1&0&-2&-1\\0&1&-1&0\end{pmatrix}.\]
Les colonnes $C_1$ et $C_2$ ne sont pas proportionnelles, $C_4=-C_1$ et $C_3=-2C_1-C_2$ : le rang vaut $2$, donc
$\dim\Ker\varphi=2$ (théorème du rang). Or $I_2$ et $A$ commutent avec $A$, et ne sont pas proportionnelles : $\Ker\varphi=\vect(I_2,A)$.
Et $\Im\varphi=\vect\big(\varphi(E_{1,1}),\varphi(E_{1,2})\big)$, un plan contenu dans l'hyperplan des matrices de trace nulle.
\q{3.} $\Ker\varphi$ et $\Im\varphi$ sont stables par $\varphi$ (proposition~1.17 du livret). Une matrice de $\Ker\varphi\cap\Im\varphi$
s'écrit $aI_2+bA$ avec une trace nulle : $2a=0$ ($\tr A=0$), donc $bA\in\Im\varphi$. Or $A=\begin{pmatrix}1&1\\1&-1\end{pmatrix}$ n'est
pas dans $\Im\varphi$ : si $A=\lambda\varphi(E_{1,1})+\mu\varphi(E_{1,2})=\begin{pmatrix}-\mu&-\lambda+2\mu\\\lambda&\mu\end{pmatrix}$,
le coefficient $(2,1)$ donne $\lambda=1$, le coefficient $(1,1)$ donne $\mu=-1$, et le coefficient $(1,2)$ vaudrait $-3\neq1$. Donc $b=0$ : l'intersection est
nulle, et $\dim\Ker\varphi+\dim\Im\varphi=4$ : ils sont supplémentaires.
\retenir{$M\mapsto AM-MA$ : son noyau contient toujours $\vect(I_n,A)$, son image est dans l'hyperplan des matrices de trace nulle.}
\par\hfill\textcolor{BleuClair}{$\square$}\par
\correction{17}{Matrices nilpotentes de taille 2}{aln-11}{énoncé : TD, p.~7}
\q{1.} Avec $M=\begin{pmatrix}a&b\\c&d\end{pmatrix}$ : $M^2=\begin{pmatrix}a^2+bc&b(a+d)\\c(a+d)&d^2+bc\end{pmatrix}$, et
$M^2-(a+d)M=\begin{pmatrix}bc-ad&0\\0&bc-ad\end{pmatrix}=-\det(M)I_2$.
\q{2.} Si $\tr A=\det A=0$, la question 1 donne $A^2=0$. Réciproquement, si $A^k=0$ : $(\det A)^k=\det(A^k)=0$ (déterminant d'un
produit de matrices $2\times2$), donc $\det A=0$, puis $A^2=\tr(A)A$ et, par récurrence, $A^k=\tr(A)^{k-1}A$. Si $A\neq0$, cela impose
$\tr A=0$ ; si $A=0$, c'est clair.
\q{3.} $\mathcal N$ est contenu dans l'hyperplan $\Ker\tr$, de dimension $3$, qui est un sous-espace : $\vect(\mathcal N)\subset\Ker\tr$.
Réciproquement, $E_{1,2}$, $E_{2,1}$ et $\begin{pmatrix}1&1\\-1&-1\end{pmatrix}$ sont nilpotentes (trace et déterminant nuls) et
engendrent $\Ker\tr$ : la troisième moins la première plus la deuxième donne $\mathrm{diag}(1,-1)$.
\retenir{en taille $2$, $M^2-\tr(M)M+\det(M)I_2=0$ : un calcul direct, très utile.}
\par\hfill\textcolor{BleuClair}{$\square$}\par
\correction{18}{$AB-BA=I$ en dimension infinie}{aln-13}{énoncé : TD, p.~7}
\q{1.} $(D\circ M)(P)=(XP)'=P+XP'$ et $(M\circ D)(P)=XP'$ : la différence vaut $P$.
\q{2.} L'argument de la trace ($\tr(AB-BA)=0\neq n$) n'a de sens qu'en dimension finie. Ici $E$ est de dimension infinie, et $M$ ne
laisse stable aucun $\R_n[X]$ (elle augmente le degré) : on ne peut pas se ramener à des matrices.
\retenir{les résultats sur la trace supposent la dimension finie.}
\par\hfill\textcolor{BleuClair}{$\square$}\par
\partie{IV}{Exercices difficiles et type concours}{}
\correction{19}{Endomorphismes de $\mathbb C^2$ de trace nulle}{aln-10}{énoncé : TD, p.~8}
\choix{la relation $M^2-\tr(M)M+\det(M)I_2=0$ (exercice sur les matrices nilpotentes) ramène tout à $f\circ f=\lambda\Id$.}
\q{1.} Avec $M$ la matrice de $f$ dans une base : $\tr M=0$ donne $M^2=-\det(M)I_2$, soit $f\circ f=\lambda\Id$ avec $\lambda=-\det f$.
\q{2.} Soit $\alpha\in\C$ tel que $\alpha^2=\lambda$ ($\alpha\neq0$) et $s=\frac1\alpha f$ : $s\circ s=\frac{\lambda}{\alpha^2}\Id=\Id$, donc
$s$ est une symétrie et $f=\alpha s$. Comme $\tr s=0$, $s\neq\pm\Id$ (de traces $\pm2$) : $\Ker(s-\Id)=\vect(e_1)$ et
$\Ker(s+\Id)=\vect(e_2)$ sont deux droites supplémentaires. Dans la base $(e_1+e_2,e_1-e_2)$, $s$ échange les deux vecteurs, sa matrice
est $\begin{pmatrix}0&1\\1&0\end{pmatrix}$, et celle de $f$ est $\begin{pmatrix}0&\alpha\\\alpha&0\end{pmatrix}$.
\q{3.} $f\circ f=0$ donne $\Im f\subset\Ker f$ ; comme $f\neq0$, $\rg f\geq1$, et le théorème du rang impose $\rg f=\dim\Ker f=1$, donc
$\Im f=\Ker f$. Soit $v\notin\Ker f$ : $\big(f(v),v\big)$ est libre ($f(v)\in\Ker f$, $v\notin\Ker f$), c'est une base, et
$f\big(f(v)\big)=0$, $f(v)=1\cdot f(v)$ : la matrice est $\begin{pmatrix}0&1\\0&0\end{pmatrix}$.
\retenir{en dimension $2$, une trace nulle donne $f\circ f=\lambda\Id$ : symétrie à un facteur près, ou carré nul.}
\par\hfill\textcolor{BleuClair}{$\square$}\par
\end{document}
